\documentclass[10pt]{article} 
\renewcommand{\thempfootnote}{\arabic{mpfootnote}}

\usepackage[hidelinks]{hyperref}

\usepackage[utf8]{inputenc}

\usepackage[T1]{fontenc}
\usepackage[ttdefault=true]{AnonymousPro}

\usepackage{multicol}
\setlength{\columnsep}{1cm}

\usepackage{abc}
\usepackage{graphicx}
\usepackage{musicography}

\usepackage[margin=0.5in]{geometry} 
\geometry{letterpaper}

\usepackage{fancyhdr} 
\pagestyle{fancy} % options: empty , plain , fancy
\renewcommand{\headrulewidth}{0pt} % customise the layout...
\lhead{}\chead{}\rhead{}
\lfoot{}\cfoot{\thepage}\rfoot{}

\begin{document}

\thispagestyle{empty}

\begin{abc}[name=shadygrove,options={-O= +c -F single.fmt}]
X:1
T:Shady Grove
M:2/4
L:1/8
K:G    
A>A A2 | BAG2 | A>A cd | e2e2 | g>g e2 | dB/A/ G2 | AB dB | A2 A2 |]
\end{abc}


\noindent\makebox[\textwidth][c]{
  \begin{minipage}{7.0in}
    \setlength{\parindent}{1.5em}


The tune Shady Grove is associated with the Upper South area in the
Appalachian Mountains having been revitalized by Jean Ritchie in the
1950s and Doc Watson in the 1960s (among others). Sources claim the
tune dates back to 1600s British Isles with its first American version
published in the early 20th
century.\footnote{\url{https://alexbledsoe.com/2010/02/08/shady-grove-and-the-tradition-of-living-songs/};
  \url{https://www.loc.gov/item/afcreed000106/}} Shady Grove is in the
same tune family as Boyne Water, Barbara Allen, and Matty Groves each
being so old that it is unclear which tune came
first.\footnote{\textit{ibid.}} The tune's simple strophic form and
dark Dorian mode have inspired at least 300 stanzas easily allowing
the performer to change the song from an innocent lovesick ballad to a
dark cautionary tale.\footnote{Ritchie, Fiona and Doug
  Orr. \textbf{Wayfaring Strangers: The Musical Heritage from Scotland
    and Ulster to Appalachia}, 2014, p xv.} \footnote{I always thought
  that `Shady Grove' was a symbol of the afterlife, much like `going
  home' and `trains' are symbols of dying. I did not find this
  mentioned while doing research but I cannot be the only one who came
  up with that, right?}

\begin{multicols}{2}
      {\noindent
      \raggedright
      \small{
      Chorus: \\
Shady grove, my true love, \\
Shady grove, I know, \\
Shady grove, my true love, \\
I'm bound for the shady grove. \\
~\\
Verse: \\
Peaches in the summertime, \\
Apples in the fall, \\
If I can't get the girl I love, \\
Won't have none at all. \\
      }
      }
      \columnbreak
          {\noindent
            \small{
When I wus a little boy, \\
I wanted a whittlin' knife; \\
Now I am a great big boy \\
An' I want a little wife... \\
~\\
Wish I had a banjo string, \\
Made of golden twine, \\
And every tune I'd pick on it – \\
Is `I wish that girl were mine.'}\footnote{Lomax, Alan. \textbf{Folk Songs of North America},
  New York, 1960, p 234.} \\
}
    \end{multicols}
 
---Kitty Steetle


  \end{minipage}
}

\begin{abc}[name=shadygrove2,options={-O= +c -F single.fmt}]
X:1
T:Shady Grove (major version)
M:C|
L:1/8
K:A  
AA | "E"B2 B2 B2 B2 | "A"cABc A2A2 | "E"B2 B2 "A"c2 e2 | "D"f4 f4 | 
"D"ef-fg aefa | "A"efed c2 c2- | c2 c2 ec-c2 | "E"B4 B2 :| 
fg | "D"a2 fg a2f2 | "A"efed c2 fg | "D"a2a2 aa-a2 | f4 f4 | 
"D"ef-fg aefa | "A"efed c2 c2- | c2 c2 ec-c2 | "E"B4 B2 :| 
\end{abc}

\noindent\makebox[\textwidth][c]{
  \begin{minipage}{7.0in}
    \setlength{\parindent}{1.5em}
    There's another ``Shady Grove'' that I like, similar to the above but in a more major mode. It apparently can be traced back to autoharp player Kilby Snow, but I think I first heard it from the Chicken Chokers. Fun fact: you can sing verses from ``June Apple'', ``Fly Around My Pretty Little Miss'', ``The Blackest Crow'', and many other songs to this melody (and vice versa), and people do.

    ---Josh Larios
  \end{minipage}
}

\end{document}
